% Book 4, Year 2, chapter 3 -- Inheritance and polymorphism.
% Provenance (never printed): sources/book4/BRIEFS.md, chapter 3.
\chapter{Inheritance and Polymorphism}\label{ch:b2:cpp-polymorphism}

The town library fines late borrowers, and not every item is fined alike. Ordinary books cost 10 cents a day, at
most 5.00; comics, always in demand, 25 cents a day, at most 3.00; study books get three days of grace, then 20
cents a day with no ceiling. In C the natural program gives each item a kind tag and switches on it.

\omcode{code/bachelor-2/03-cpp-polymorphism/listings/fines_switch.c}{4}{34}{Fines by a switch on a kind
tag}\label{lst:b2:cpp-polymorphism:fines-switch}

\omoutput{code/bachelor-2/03-cpp-polymorphism/outputs/fines_switch.txt}{Three items, three rules, three
lateness values.}\label{out:b2:cpp-polymorphism:fines-switch}

The program is correct, and it does not age well. Every function that depends on the kind repeats the switch: here
two of them, in a real system dozens (the loan period, the shelf, the label, the reminder letter). When the library
starts lending maps, someone must find every switch and add a case to each, and the compiler warns only for some
of the switches it can see. This chapter gives each kind its own class and lets the call itself choose the code of
the right kind: adding a kind then means adding a class, in a new file, and editing nothing that works.

\begin{recall}
A class (\cref{ch:b2:cpp-classes-raii}) bundles data with the functions that use them, with an invariant kept by
its \omterm{def:b2:cpp-classes-raii:constructor}{constructor} and members. The school volumes met \omterm{def:b2:cpp-polymorphism:inheritance}{inheritance}, in Python: a class can be declared as a special
case of another and inherit its methods. In C, a pointer to a function (Book 3) holds the address of a function,
and a call through it runs whichever function it points to: the choice is made at run time. Objects are copied and
moved by the rules of \cref{ch:b2:cpp-copy-move}.
\end{recall}

\section{Derived classes}\label{sec:b2:cpp-polymorphism:derived}

\begin{definition}[Base class, derived class, inheritance]\label{def:b2:cpp-polymorphism:inheritance}
A class \code{D} declared as \code{class D : public B} is a \emph{derived class}\index{derived class} of the
\emph{base class}\index{base class} \code{B}. Every \code{D} object contains a \code{B} object (its \emph{base
part}) followed by the members \code{D} declares, and \code{D} has every \omterm{def:b2:cpp-classes-raii:class}{member function} of \code{B}. This is
\emph{inheritance}\index{inheritance}. A \code{D} can be used wherever a \code{B} is expected: a \code{D *} converts
to a \code{B *} and a \code{D} binds to a \code{B &}, both pointing at the base part.
\end{definition}

Access gains a third level. A \code{public} member is visible to everyone, a \code{private} one only to the class's
own members, and a \code{protected} one to the class's members and to those of the classes derived from it. Since
a derived object contains a base object, building one runs two \omterm{def:b2:cpp-classes-raii:constructor}{constructors}.

\omcode{code/bachelor-2/03-cpp-polymorphism/listings/ctor_order.cpp}{5}{32}{A \code{Map} is an
\code{Item}}\label{lst:b2:cpp-polymorphism:ctor-order}

\omoutput{code/bachelor-2/03-cpp-polymorphism/outputs/ctor_order.txt}{Base first on the way in, last on the way
out.}\label{out:b2:cpp-polymorphism:ctor-order}

The order is fixed by the language: the base part is constructed first (its members, then its \omterm{def:b2:cpp-classes-raii:constructor}{constructor}'s body),
then the \omterm{def:b2:cpp-polymorphism:inheritance}{derived class}'s members in order of declaration, then its body; destruction runs the same steps in
reverse. When a derived \omterm{def:b2:cpp-classes-raii:constructor}{constructor}'s body starts, its base part is complete and may be used; when a base
\omterm{def:b2:cpp-classes-raii:constructor}{destructor}'s body runs, the derived part is already gone. A derived \omterm{def:b2:cpp-classes-raii:constructor}{constructor} passes arguments to the base
\omterm{def:b2:cpp-classes-raii:constructor}{constructor} in its initialiser list, as \code{Recording(...) : Item(std::move(t))} does below.

The rule ``a \code{D} can be used where a \code{B} is expected'' has a trap when the \code{B} is an object rather
than a reference or a pointer. A \code{B} object has room for a \code{B} and nothing more.

\begin{definition}[Object slicing]\label{def:b2:cpp-polymorphism:slicing}
\emph{Object slicing}\index{object slicing} is the copy of a derived object into an object of its \omterm{def:b2:cpp-polymorphism:inheritance}{base class}: by
assignment, by initialisation, or by passing it by value to a parameter of the base type. Only the base part is
copied, and the result is a plain base object; the derived members and the derived behaviour are lost.
\end{definition}

\omcode{code/bachelor-2/03-cpp-polymorphism/listings/slicing.cpp}{5}{38}{A recording passed by reference and by
value}\label{lst:b2:cpp-polymorphism:slicing}

\omoutput{code/bachelor-2/03-cpp-polymorphism/outputs/slicing.txt}{By value, the recording lost its
minutes.}\label{out:b2:cpp-polymorphism:slicing}

\begin{omfigure}
\begin{tikzpicture}[font=\footnotesize]
\node[font=\footnotesize\bfseries, anchor=west] at (-0.2,2.0) {\texttt{Recording r} (48 bytes)};
\node[cell, minimum width=1.6cm, fill=omDef!10] (t) at (0.8,1.2) {table ptr};
\node[cell, minimum width=2.6cm, fill=omDef!10, right=0mm of t] (ti) {\texttt{title\_}};
\node[cell, minimum width=1.4cm, fill=omThm!10, right=0mm of ti] (mi) {\texttt{minutes\_}};
\draw[decorate, decoration={brace, mirror, amplitude=4pt}] (t.south west) -- (ti.south east)
  node[midway, below=4pt, font=\scriptsize] {base part: an \texttt{Item}};
\node[font=\footnotesize\bfseries, anchor=west] at (7.4,2.0) {\texttt{Item copy = r;} (40 bytes)};
\node[cell, minimum width=1.6cm, fill=omDef!10] (t2) at (8.2,1.2) {table ptr};
\node[cell, minimum width=2.6cm, fill=omDef!10, right=0mm of t2] (ti2) {\texttt{title\_}};
\draw[->, thick, omThm] (ti.north) -- ++(0,0.25) -| (ti2.north) node[pos=0.25, above, font=\scriptsize] {copy};
\node[font=\scriptsize, text=omThm, align=center, below=1mm of mi] {\texttt{minutes\_} has no\\ room: it is cut off};
\node[font=\scriptsize, text=black!60, align=center] at (9.5,0.2) {the table pointer is \texttt{Item}'s:\\ \texttt{label()} is \texttt{Item::label}};
\end{tikzpicture}
\par\smallskip{\small Slicing. Copying a \code{Recording} into an \code{Item} copies the base part only. The copy is
a genuine \code{Item}, whose \omterm{def:b2:cpp-polymorphism:vtable}{table pointer} (\cref{sec:b2:cpp-polymorphism:virtual}) is that of \code{Item}, so even its
virtual calls run \code{Item}'s functions.}
\end{omfigure}

\section{Virtual functions}\label{sec:b2:cpp-polymorphism:virtual}

\begin{definition}[Virtual function, override]\label{def:b2:cpp-polymorphism:virtual}
A \omterm{def:b2:cpp-classes-raii:class}{member function} declared \code{virtual} in a \omterm{def:b2:cpp-polymorphism:inheritance}{base class} is a \emph{virtual function}\index{virtual function}. A
\omterm{def:b2:cpp-polymorphism:inheritance}{derived class} \emph{overrides}\index{override} it by declaring a function with the same name and parameters; the
keyword \code{override} asks the compiler to check that there is something to override, and \code{final} forbids
further overriding. A function that is not virtual can only be \emph{hidden} by a derived function of the same name,
which is a different thing.
\end{definition}

\begin{definition}[Dynamic dispatch, subtype polymorphism]\label{def:b2:cpp-polymorphism:dispatch}
A call of a \omterm{def:b2:cpp-polymorphism:virtual}{virtual function} through a pointer or a reference runs the version of the class of the object pointed
at, which is known only at run time, not the version of the pointer's type: this is \emph{dynamic
dispatch}\index{dynamic dispatch}. Code written against a \omterm{def:b2:cpp-polymorphism:inheritance}{base class} and able to work, through dynamic dispatch, with
objects of every \omterm{def:b2:cpp-polymorphism:inheritance}{derived class} has \emph{subtype polymorphism}\index{subtype polymorphism}.
\end{definition}

The shapes below share one \omterm{def:b2:cpp-polymorphism:inheritance}{base class}. \code{area} has no sensible default and is declared \code{= 0} (the next
section explains); \code{name} has one; \code{describe} is an ordinary function written once in the \omterm{def:b2:cpp-polymorphism:inheritance}{base class},
whose calls to \code{name} and \code{area} are dispatched.

\omcode{code/bachelor-2/03-cpp-polymorphism/listings/shapes.hpp}{5}{31}{A \omterm{def:b2:cpp-polymorphism:inheritance}{base class} and two derived
classes}\label{lst:b2:cpp-polymorphism:shapes}

\omcode{code/bachelor-2/03-cpp-polymorphism/listings/shapes.cpp}{6}{19}{A plan of shapes, each called through
\code{Shape}}\label{lst:b2:cpp-polymorphism:shapes-main}

\omoutput{code/bachelor-2/03-cpp-polymorphism/outputs/shapes.txt}{Each call ran the function of the real
class.}\label{out:b2:cpp-polymorphism:shapes}

The loop sees only \code{const Shape &}, yet each \code{describe} printed the right name and area. Destroying the
objects through the \omterm{def:b2:cpp-polymorphism:inheritance}{base class} needs one more \omterm{def:b2:cpp-polymorphism:virtual}{virtual function}. When \code{delete} is applied to a \code{B *} that
points at a \code{D}, the \omterm{def:b2:cpp-classes-raii:constructor}{destructor} called is \code{B}'s unless \code{B}'s \omterm{def:b2:cpp-classes-raii:constructor}{destructor} is virtual, and the behaviour
is undefined.

\omcode{code/bachelor-2/03-cpp-polymorphism/listings/no_virtual_dtor.cpp}{6}{22}{A derived object deleted
through its base}\label{lst:b2:cpp-polymorphism:no-virtual-dtor}

\omoutput[1-7,18-18]{code/bachelor-2/03-cpp-polymorphism/outputs/no_virtual_dtor.txt}{AddressSanitizer: the size
freed is not the size allocated.}\label{out:b2:cpp-polymorphism:no-virtual-dtor}

Only \code{~Item} ran: the four strings of the recording were never destroyed, and \code{delete} returned the block
as if it held a 4-byte \code{Item} when 136 bytes had been allocated. The rule follows: a class meant to be used
through base pointers declares \code{virtual ~B() = default;}, as \code{Shape} does. A \code{std::unique_ptr<B>} does
not help, since it calls the same \code{delete} (\cref{exo:b2:cpp-polymorphism:2}).

How does a call find the right function at run time? The compiler gives every class with \omterm{def:b2:cpp-polymorphism:virtual}{virtual functions} one
table of function addresses, and every object of such a class one hidden pointer to the table of its class, set by
the \omterm{def:b2:cpp-classes-raii:constructor}{constructor}. The same mechanism can be written by hand in C.

\begin{definition}[Dispatch table, table pointer]\label{def:b2:cpp-polymorphism:vtable}
The \emph{dispatch table}\index{dispatch table} (or \emph{vtable}\index{vtable}) of a class with \omterm{def:b2:cpp-polymorphism:virtual}{virtual functions} is
an array, one per class and built by the compiler, of the addresses of the versions of its \omterm{def:b2:cpp-polymorphism:virtual}{virtual functions} that
apply to that class. Each object of the class holds a \emph{table pointer}\index{table pointer} to its class's table.
A virtual call loads the table pointer from the object, loads the address from a fixed position in the table, and
calls it.
\end{definition}

\omcode{code/bachelor-2/03-cpp-polymorphism/listings/vtable_by_hand.c}{4}{41}{\omterm{def:b2:cpp-polymorphism:vtable}{Dispatch tables} written in
C}\label{lst:b2:cpp-polymorphism:vtable-c}

\omoutput{code/bachelor-2/03-cpp-polymorphism/outputs/vtable_by_hand.txt}{The first two lines of the
output of \cref{lst:b2:cpp-polymorphism:shapes-main}, by hand.}\label{out:b2:cpp-polymorphism:vtable-c}

\begin{omfigure}
\begin{tikzpicture}[font=\footnotesize]
\foreach \nm/\y/\cls/\col in {bed/2.4/rect/omDef, pond/0/circle/omThm} {
  \node[anchor=east] at (-0.1,\y) {\texttt{\nm}};
  \node[cell, minimum width=1.2cm, fill=\col!10] (\nm p) at (0.6,\y) {$\bullet$};
  \node[cell, minimum width=1.9cm, right=0mm of \nm p, font=\scriptsize] {data of a \cls};
}
\node[draw, fill=omDef!10, align=left, minimum width=3.0cm, font=\scriptsize\ttfamily] (RT) at (5.6,2.4)
  {RECT\_TABLE\\ {[0]} rect\_area\\ {[1]} rect\_name};
\node[draw, fill=omThm!10, align=left, minimum width=3.0cm, font=\scriptsize\ttfamily] (CT) at (5.6,0.0)
  {CIRCLE\_TABLE\\ {[0]} circle\_area\\ {[1]} circle\_name};
\draw[ptr] (bedp.center) -- (0.6,3.0) -- (3.4,3.0) -- (RT.west);
\draw[ptr] (pondp.center) -- (0.6,-0.6) -- (3.4,-0.6) -- (CT.west);
\node[draw, rounded corners=2pt, align=left, font=\scriptsize\ttfamily, fill=black!4] (code) at (10.2,1.2)
  {rect\_area: w*h\\ rect\_name: its name\\ circle\_area: pi*r*r\\ circle\_name: its name};
\draw[ptr, black!60] (RT.east) -- (code.west);
\draw[ptr, black!60] (CT.east) -- (code.west);
\node[font=\scriptsize, text=black!60, align=center] at (1.3,-1.2) {objects: one table pointer each};
\node[font=\scriptsize, text=black!60, align=center] at (5.6,-1.2) {tables: one per class};
\node[font=\scriptsize, text=black!60, align=center] at (10.2,-1.2) {code: one copy of each function};
\end{tikzpicture}
\par\smallskip{\small The two objects of \cref{lst:b2:cpp-polymorphism:vtable-c}, of two classes. A call \code{s->table->area(s)}, or its C++ form
\code{s->area()}, follows the object's \omterm{def:b2:cpp-polymorphism:vtable}{table pointer} to its class's table, takes the address at the position
reserved for \code{area}, and calls that function with the object.}
\end{omfigure}

\begin{proposition}[Cost of dispatch tables]\label{prop:b2:cpp-polymorphism:layout}
In this scheme, a program with $k$ classes that have \omterm{def:b2:cpp-polymorphism:virtual}{virtual functions} holds $k$ tables, the table of a class
having one entry per \omterm{def:b2:cpp-polymorphism:virtual}{virtual function} of the class, and each object of these classes holds exactly one \omterm{def:b2:cpp-polymorphism:vtable}{table
pointer}, whatever the number of \omterm{def:b2:cpp-polymorphism:virtual}{virtual functions}. A virtual call costs two loads (the \omterm{def:b2:cpp-polymorphism:vtable}{table pointer} and the
address) and an indirect call, independently of the number of classes.
\end{proposition}

\begin{proof}
The tables depend on the class only: all objects of one class call the same versions, so one table per class
suffices, with one slot per \omterm{def:b2:cpp-polymorphism:virtual}{virtual function}, a \omterm{def:b2:cpp-polymorphism:inheritance}{derived class}'s table extending its base's with the same positions
for the inherited functions so that a call compiled against the base finds the right slot. An object needs only to
say which class it is of, and one pointer to that class's table says it, whatever the table's length. The call
sequence of \cref{def:b2:cpp-polymorphism:vtable} reads the object's pointer and one slot at a position fixed at
compile time, then calls; nothing in it depends on how many classes exist.
\end{proof}

\omcode{code/bachelor-2/03-cpp-polymorphism/listings/sizes.cpp}{4}{29}{Sizes of objects with and without \omterm{def:b2:cpp-polymorphism:virtual}{virtual
functions}}\label{lst:b2:cpp-polymorphism:sizes}

\omoutput{code/bachelor-2/03-cpp-polymorphism/outputs/sizes.txt}{One \omterm{def:b2:cpp-polymorphism:vtable}{table pointer}, whatever the number of \omterm{def:b2:cpp-polymorphism:virtual}{virtual
functions}.}\label{out:b2:cpp-polymorphism:sizes}

The C++ standard does not prescribe tables, but every major compiler uses them, with the \omterm{def:b2:cpp-polymorphism:vtable}{table pointer} at the start
of the object, as the last line of the output shows. On this 64-bit machine a pointer is 8 bytes: \code{Plain} is
its \code{int}; \code{OneVirtual} and \code{ManyVirtual} are a pointer plus an \code{int}, padded to a multiple of 8;
and \code{Derived}'s extra \code{int} fits in that padding (\cref{exo:b2:cpp-polymorphism:4}).

\section{Abstract classes}\label{sec:b2:cpp-polymorphism:abstract}

\begin{definition}[Pure virtual function, abstract class]\label{def:b2:cpp-polymorphism:abstract}
A \omterm{def:b2:cpp-polymorphism:virtual}{virtual function} declared with \code{= 0}, as \code{virtual double area() const = 0;}, is a \emph{pure virtual
function}\index{pure virtual function}: it has no version in its class, and every \omterm{def:b2:cpp-polymorphism:inheritance}{derived class} that is to be
instantiated must \omterm{def:b2:cpp-polymorphism:virtual}{override} it. A class with at least one pure \omterm{def:b2:cpp-polymorphism:virtual}{virtual function}, declared or inherited and not
overridden, is an \emph{abstract class}\index{abstract class}: no object of it can be created, only objects of
derived classes that \omterm{def:b2:cpp-polymorphism:virtual}{override} every pure function.
\end{definition}

An \omterm{def:b2:cpp-polymorphism:abstract}{abstract class} with only \omterm{def:b2:cpp-polymorphism:abstract}{pure virtual functions} and a virtual \omterm{def:b2:cpp-classes-raii:constructor}{destructor} describes what its derived classes
provide and nothing of how. It plays the role an interface plays in other languages, and code written against it
can be compiled before any of its implementations exists. \code{Shape} is such a class but for \code{describe} and
\code{name}, which every shape shares.

One question remains: which version runs when a \omterm{def:b2:cpp-classes-raii:constructor}{constructor} calls a \omterm{def:b2:cpp-polymorphism:virtual}{virtual function}? While the base \omterm{def:b2:cpp-classes-raii:constructor}{constructor}
runs, the derived members have not been built yet (\cref{out:b2:cpp-polymorphism:ctor-order}); calling a derived
version would read them uninitialised. C++ therefore treats the object, during the base \omterm{def:b2:cpp-classes-raii:constructor}{constructor} and the base
\omterm{def:b2:cpp-classes-raii:constructor}{destructor}, as an object of the \omterm{def:b2:cpp-polymorphism:inheritance}{base class}: virtual calls made there run the base versions
(\cref{exo:b2:cpp-polymorphism:5}). In terms of tables, each \omterm{def:b2:cpp-classes-raii:constructor}{constructor} sets the \omterm{def:b2:cpp-polymorphism:vtable}{table pointer} to its own class's
table on entry, so the pointer is the base's while the base \omterm{def:b2:cpp-classes-raii:constructor}{constructor} runs.

\section{Designing with classes}\label{sec:b2:cpp-polymorphism:design}

\omterm{def:b2:cpp-polymorphism:inheritance}{Inheritance} says ``a \code{D} is a \code{B}'': every \code{D} can stand in for a \code{B}. It is often used where a
weaker relation is meant.

\begin{definition}[Composition]\label{def:b2:cpp-polymorphism:composition}
A class uses \emph{composition}\index{composition} when it holds an object of another class as a member (``has a'')
and calls it, rather than deriving from it (``is a''). A member may be held by value, or by reference or pointer to an
object of an \omterm{def:b2:cpp-polymorphism:abstract}{abstract class}, in which case the behaviour it provides can be chosen at run time.
\end{definition}

A car has an engine and is not one; a loan has a fine policy. \omterm{def:b2:cpp-polymorphism:composition}{Composition} exposes only what the class chooses to
forward, whereas a \omterm{def:b2:cpp-polymorphism:inheritance}{derived class} inherits the whole public interface of its base, along with every promise the base
makes. Those promises are what \omterm{def:b2:cpp-polymorphism:inheritance}{inheritance} must keep.

% ledger: b2:lsp-1987
% ledger: b2:lsp-1994
\begin{definition}[Substitution principle]\label{def:b2:cpp-polymorphism:lsp}
A \omterm{def:b2:cpp-polymorphism:inheritance}{derived class} \code{D} of \code{B} satisfies the \emph{substitution principle}\index{substitution principle} if
every property that a client can prove about objects of \code{B} from the specification of \code{B} (the
preconditions and postconditions of its functions, and its invariant) also holds for objects of \code{D}.
Equivalently, a correct client of \code{B} stays correct when given a \code{D}. The principle bears the name of
Barbara Liskov, who stated it in a keynote in 1987 and, with Jeannette Wing, as a definition of subtyping in 1994.
\end{definition}

Geometry says that a square is a rectangle. Here is what happens when a class \code{Square} derives from a class
\code{Rectangle} whose setters promise to change one side only.

\begin{example}[A square that is not a rectangle]\label{ex:b2:cpp-polymorphism:square}
\omcode{code/bachelor-2/03-cpp-polymorphism/listings/square.cpp}{4}{39}{A client of \code{Rectangle} given a
\code{Square}}
\omoutput{code/bachelor-2/03-cpp-polymorphism/outputs/square.txt}{The client's promise holds for one of the
two.}\label{out:b2:cpp-polymorphism:square}
\code{resize_plot} relies on the postconditions of \code{Rectangle}: after \code{set_width(4)} and
\code{set_height(5)} the area is 20. \code{Square} must keep its sides equal, so its \code{set_height(5)} also changes
the width, and the client sees 25. The square of geometry is a rectangle because a geometric rectangle never
changes; a mutable \code{Square} cannot keep the promises of a mutable \code{Rectangle}, whichever way it \omterm{def:b2:cpp-polymorphism:virtual}{overrides}
the setters (\cref{exo:b2:cpp-polymorphism:8}). Either the two classes have no setters, or \code{Square} is not
derived from \code{Rectangle}.
\end{example}

The library's fines are the case where \omterm{def:b2:cpp-polymorphism:inheritance}{inheritance} fits. A policy is a function of the days late, and every policy
can stand in for every other: a loan asks its policy for a fine and accepts any answer. The policies form an
\omterm{def:b2:cpp-polymorphism:abstract}{abstract class} and three derived ones; the loan, or here the function that sums the fines, holds a policy by
\omterm{def:b2:cpp-polymorphism:composition}{composition}. The rule that maps a book to its policy is written once, in \code{policy_for}.

\omcode{code/bachelor-2/03-cpp-polymorphism/listings/library.hpp}{65}{100}{The fine policies, and the only
function that knows the genres}\label{lst:b2:cpp-polymorphism:library}

The rest of \code{library.hpp} reads \code{books.csv} and \code{loans.csv} of the town (the book's genre, and each
loan's dates; a loan not yet returned has an empty return date) and counts the days late at the end of 2025. The
program below sums the fines.

\omcode{code/bachelor-2/03-cpp-polymorphism/listings/library_demo.cpp}{12}{27}{Fines on all the town's loans; the
call \code{p.fine_cents(d)} is dispatched}\label{lst:b2:cpp-polymorphism:library-demo}

\omoutput{code/bachelor-2/03-cpp-polymorphism/outputs/run_library.txt}{The fines owed on 31 December
2025.}\label{out:b2:cpp-polymorphism:library}

\begin{omfigure}
\begin{tikzpicture}[font=\footnotesize, cls/.style={draw, fill=omDef!8, minimum width=2.6cm, align=center, inner sep=3pt}]
\node[cls] (loan) at (0,1.6) {\textbf{fine summing}\\[1pt] \scriptsize\texttt{policy\_for(book)}};
\node[cls, fill=omThm!8] (pol) at (5.0,1.6) {\textit{\textbf{FinePolicy}}\\[1pt] \scriptsize\texttt{fine\_cents(d) = 0}\\ \scriptsize\texttt{name() = 0}};
\node[cls] (s) at (2.2,-0.4) {\textbf{StandardFine}};
\node[cls] (c) at (5.0,-0.4) {\textbf{ComicFine}};
\node[cls] (st) at (7.8,-0.4) {\textbf{StudyFine}};
\foreach \n in {s,c,st} \draw[-{Triangle[open, length=2.5mm, width=2.5mm]}] (\n.north) -- (pol.south);
\draw[-{Stealth}, thick, black!70] (loan.east) -- node[above, font=\scriptsize] {has a} (pol.west);
\node[cls, dashed, fill=white] (map) at (10.6,-0.4) {\textbf{MapFine}\\ \scriptsize (new file)};
\draw[dashed, -{Triangle[open, length=2.5mm, width=2.5mm]}] (map.north) -- (pol.south east);
\node[font=\scriptsize, text=black!60, align=left, anchor=west] at (9.0,1.6) {open arrow: ``is a'' (inherits)\\ filled arrow: ``has a'' (holds)};
\end{tikzpicture}
\par\smallskip{\small The design of the fines. The code that sums fines holds a policy (\omterm{def:b2:cpp-polymorphism:composition}{composition}) and knows only
the \omterm{def:b2:cpp-polymorphism:abstract}{abstract class}; the three policies derive from it (\omterm{def:b2:cpp-polymorphism:inheritance}{inheritance}). A fourth policy, for maps, is one more class
in one more file, plus one line in \code{policy_for}.}
\end{omfigure}

\begin{method}[Adding a kind without editing working code]\label{met:b2:cpp-polymorphism:open-closed}
To make a family of kinds open to new members:
\begin{enumerate}
\item collect what every kind must provide into an \omterm{def:b2:cpp-polymorphism:abstract}{abstract class} of \omterm{def:b2:cpp-polymorphism:abstract}{pure virtual functions}, with a virtual
\omterm{def:b2:cpp-classes-raii:constructor}{destructor};
\item write each kind as a \omterm{def:b2:cpp-polymorphism:inheritance}{derived class}, in a file of its own, overriding every pure function;
\item write the code that uses the kinds against references or pointers to the \omterm{def:b2:cpp-polymorphism:abstract}{abstract class} only, so that it
compiles without seeing any \omterm{def:b2:cpp-polymorphism:inheritance}{derived class};
\item create the objects in one place (here \code{policy_for}), the only code that names the derived classes.
\end{enumerate}
A new kind is then one new file and one new line in the creating function; no switch has to be found, and the code
of step 3 is neither edited nor recompiled. This is the \emph{open--closed} rule: open to extension, closed to
modification.
\end{method}

The price is a virtual call where the switch had a jump. The program below compares a direct call that the
compiler may not inline, a virtual call on objects of one class, the same on a random mix of two classes, and a
switch on a tag that the compiler inlined, ten million calls each.

\omcodes{code/bachelor-2/03-cpp-polymorphism/listings/dispatch_cost.cpp}{8-20,47-50}{Four ways to call an area
function (\code{ns_per} returns the time of its argument divided by \code{n})}\label{lst:b2:cpp-polymorphism:dispatch-cost}

\omoutput{code/bachelor-2/03-cpp-polymorphism/outputs/dispatch_cost.txt}{One run on the author's laptop; the digits
change from run to run.}\label{out:b2:cpp-polymorphism:dispatch-cost}

Every line costs a few nanoseconds. The virtual calls cost more than the direct one for two reasons besides the two
loads: each object sits in its own heap block reached through a \code{unique_ptr}, where the direct call walks one
contiguous array; and with mixed kinds the processor cannot guess which function comes next, a cost that
\cref{ch:b2:processor-and-pipelining} explains. For a fine computed once per loan none of this matters; in an inner
loop over millions of small objects it can, and the switch or a template (\cref{ch:b2:cpp-templates}) is then the
better tool.

\begin{history}[Simula]
% ledger: b2:simula-oop
% ledger: b2:simula-hopl
Classes, \omterm{def:b2:cpp-polymorphism:inheritance}{inheritance} and dispatch on the class of an object come from Simula 67, a language designed by Ole-Johan
Dahl and Kristen Nygaard to write simulations, whose objects were the people and machines being simulated. Its
designers told the history of the language at the first conference on the history of programming languages, in
1978. C++ was created to write systems programs in the style Simula encouraged, with the efficiency of C.
\end{history}

\section{Exercises}\label{sec:b2:cpp-polymorphism:exercises}

\begin{exercise}[$\star$]\label{exo:b2:cpp-polymorphism:1}
Predict the four lines printed. For each call, say whether it is dispatched at run time and why.
\omcode{code/bachelor-2/03-cpp-polymorphism/listings/exo_which.cpp}{4}{27}{Pointer, reference, value}
\end{exercise}

\begin{exercise}[$\star$]\label{exo:b2:cpp-polymorphism:2}
This program owns its sensors with \code{std::unique_ptr}, so nothing should go wrong; yet AddressSanitizer stops it
at the end of \code{main}. Explain the report and fix the program with one word.
\omcode{code/bachelor-2/03-cpp-polymorphism/listings/exo_owner.cpp}{7}{23}{Two sensors}
\omoutput[1-7]{code/bachelor-2/03-cpp-polymorphism/outputs/exo_owner.txt}{The start of the report.}
\end{exercise}

\begin{exercise}[$\star$]\label{exo:b2:cpp-polymorphism:3}
Add a class \code{Triangle}, given its three sides, to the shapes of \cref{lst:b2:cpp-polymorphism:shapes}, in a new
file and without editing \code{shapes.hpp} or \code{shapes.cpp}'s loop. Its \omterm{def:b2:cpp-classes-raii:constructor}{constructor} must refuse three lengths
that are not the sides of a (non-flat) triangle. Which step of \cref{met:b2:cpp-polymorphism:open-closed} does each
part of your file follow?
\end{exercise}

\begin{exercise}[$\star$]\label{exo:b2:cpp-polymorphism:4}
On this machine \code{int} is 4 bytes and aligned on 4, a pointer 8 bytes and aligned on 8, and an object's size is
a multiple of its strictest member alignment. Explain each size of \cref{out:b2:cpp-polymorphism:sizes}. What would
\code{sizeof} print for a class with a \omterm{def:b2:cpp-polymorphism:vtable}{table pointer} and three \code{char} members? For a class with two
\code{double} members and five \omterm{def:b2:cpp-polymorphism:virtual}{virtual functions}?
\end{exercise}

\begin{exercise}[$\star\star$]\label{exo:b2:cpp-polymorphism:5}
Predict the output, line by line. Why does the \omterm{def:b2:cpp-classes-raii:constructor}{constructor} of \code{Item} not print \code{a map of the coast}? Which
would be worse if it did?
\omcode{code/bachelor-2/03-cpp-polymorphism/listings/exo_ctor_virtual.cpp}{5}{22}{A virtual call in a \omterm{def:b2:cpp-classes-raii:constructor}{constructor} and
a \omterm{def:b2:cpp-classes-raii:constructor}{destructor}}
\end{exercise}

\begin{exercise}[$\star\star$]\label{exo:b2:cpp-polymorphism:6}
The catalogue below prints \code{Harbour Songs} without its minutes. Explain why, then fix it so that the catalogue
holds books and recordings and prints each with its own label. Which owner type did you choose for the elements,
and why not \code{std::vector<Item &>}?
\omcode{code/bachelor-2/03-cpp-polymorphism/listings/exo_slice_vector.cpp}{6}{30}{A catalogue}
\omoutput{code/bachelor-2/03-cpp-polymorphism/outputs/exo_slice_vector.txt}{The output.}
\end{exercise}

\begin{exercise}[$\star\star$]\label{exo:b2:cpp-polymorphism:7}
Rewrite \code{fines_switch.c} in C without any \code{switch}, with one table per kind as in
\cref{lst:b2:cpp-polymorphism:vtable-c}: a \code{struct kind_table} holding the kind's name and a pointer to its fine
function, and items that point to their kind's table. Test that your fines equal those of the switch for every
kind and every lateness from $-2$ to 400 days. What must be done to add a kind?
\end{exercise}

\begin{exercise}[$\star\star$]\label{exo:b2:cpp-polymorphism:8}
State the postconditions of \code{Rectangle::set_width} and \code{set_height}. Prove that no \omterm{def:b2:cpp-polymorphism:virtual}{override} of the two
setters in a class \code{Square} derived from \code{Rectangle} can both keep the invariant ``width equals height''
and satisfy the \omterm{def:b2:cpp-polymorphism:lsp}{substitution principle} for the client \code{resize_plot}. Then show that an immutable
\code{Rectangle} (\omterm{def:b2:cpp-classes-raii:constructor}{constructor} and \code{area} only) can have an immutable \code{Square} as a \omterm{def:b2:cpp-polymorphism:inheritance}{derived class} that
satisfies the principle.
\end{exercise}

\begin{exercise}[$\star\star$]\label{exo:b2:cpp-polymorphism:9}
The council caps every fine at 10.00. Write a policy \code{CappedFine} that holds a reference to another policy and
a cap, and returns the smaller of the two. Is it \omterm{def:b2:cpp-polymorphism:composition}{composition} or \omterm{def:b2:cpp-polymorphism:inheritance}{inheritance}, and why is it both? Compute the total
owed on the town's loans at the end of 2025 with every policy capped at 10.00, and check that a very large cap
changes nothing.
\end{exercise}

\begin{exercise}[$\star\star\star$]\label{exo:b2:cpp-polymorphism:10}
Represent arithmetic expressions with an \omterm{def:b2:cpp-polymorphism:abstract}{abstract class} \code{Expr} and derived classes for numbers, variables,
sums and products, owned through \code{std::unique_ptr}. Give them \code{eval} (in an environment mapping variables
to values), \code{str} (fully parenthesised) and \code{simplify}, which folds operations on two numbers and applies
$0 + e = e$, $1 \times e = e$ and $0 \times e = 0$. Put what sums and products share in a common \omterm{def:b2:cpp-polymorphism:inheritance}{base class}. Test
that \code{((x * (2 + 3)) + (0 * y))} simplifies to \code{(x * 5)} and evaluates as before.
\end{exercise}

\begin{exercise}[$\star\star\star$]\label{exo:b2:cpp-polymorphism:11}
A program has $k$ classes with \omterm{def:b2:cpp-polymorphism:virtual}{virtual functions}; class $i$ has $m_i$ \omterm{def:b2:cpp-polymorphism:virtual}{virtual functions} and $N_i$ objects. Count the
bytes of tables and of \omterm{def:b2:cpp-polymorphism:vtable}{table pointers}, with 8-byte pointers and one slot per \omterm{def:b2:cpp-polymorphism:virtual}{virtual function}. For the library with
3 policies of 3 \omterm{def:b2:cpp-polymorphism:virtual}{virtual functions} each (the \omterm{def:b2:cpp-classes-raii:constructor}{destructor} counts once), one object each, and 400 loans that each hold a
policy pointer, compare with the C design of \cref{lst:b2:cpp-polymorphism:fines-switch}. Then take a game with
10 kinds of particles, 4 \omterm{def:b2:cpp-polymorphism:virtual}{virtual functions} each, and $10^6$ particles: what fraction of the memory is \omterm{def:b2:cpp-polymorphism:vtable}{table
pointers} if a particle holds 24 bytes of data?
\end{exercise}

\begin{exercise}[$\star\star\star$]\label{exo:b2:cpp-polymorphism:12}
A virtual call chooses a function by the class of one object. Collision tests between two bodies, discs and boxes
with sides parallel to the axes, depend on both. Implement \code{bool hits(const Body &other)} by \emph{double
dispatch}: \code{a.hits(b)} makes a virtual call on \code{a}, which makes a second virtual call on \code{b} naming
its own class. Test all four pairs. How many functions does each class need with $c$ classes of bodies, and what
must change when a class is added?
\end{exercise}

\section{Problem: The town library's lending rules}\label{sec:b2:cpp-polymorphism:problem}

\begin{problem}[{Weekend problem --- The town library's lending rules}]\label{pb:b2:cpp-polymorphism:1}
The library of Brenmoor wants its software to handle every rule that depends on the kind of item: the fine, the loan
period and the label printed on reminders. It starts from the C program of \cref{lst:b2:cpp-polymorphism:fines-switch}
and the class design of \cref{lst:b2:cpp-polymorphism:library}, and plans to lend maps from next year, fined 50 cents
a day, at most 4.00.

\needspace{4\baselineskip}\medskip\noindent\textbf{Part I --- The rules as a switch.}\par\nopagebreak
\begin{enumerate}[beginpenalty=10000]
\item In the C program, list every place that must change to add the kind \code{MAP}.
\item The library adds two more functions on the kind: the loan period (21 days for every kind but study books, 28)
and the shelf. With $f$ such functions and $k$ kinds, how many \code{case}s does the C program contain? How many must
be written to add one kind, and in how many places?
\item With \code{-Wall}, \code{gcc} warns about a \code{switch} on an \code{enum} that misses a value, but only if
the switch has no \code{default}. Which switches of the program would warn after \code{MAP} is added?
\item Using the fines of \cref{out:b2:cpp-polymorphism:fines-switch}, give the fine of a map 2, 10 and 60 days late.
\item Why can a C program read in from files (the kind as a string in a catalogue) not escape a switch or a table
somewhere, even with classes?
\item In which file of the class design does that unavoidable choice live?
\end{enumerate}

\needspace{4\baselineskip}\medskip\noindent\textbf{Part II --- A class hierarchy.}\par\nopagebreak
\begin{enumerate}[resume, beginpenalty=10000]
\item Write the class \code{MapFine}. Which files change, and which do not?
\item Add the loan period and the shelf to \code{FinePolicy}. Should the class keep its name? Propose one, and say
which method of the chapter you follow.
\item With $f$ \omterm{def:b2:cpp-polymorphism:virtual}{virtual functions} and $k$ kinds, how many function bodies does the design contain? Compare with
question 2.
\item From \cref{out:b2:cpp-polymorphism:library}, give the number of late loans and the mean fine per late loan for
each policy.
\item Loan 27 owes 62.40 under the study policy. How many days late is it? Was it returned?
\item The library council finds the study policy unfair to long loans and caps it at 10.00, using the
\code{CappedFine} of \cref{exo:b2:cpp-polymorphism:9}. Which file changes?
\end{enumerate}

\needspace{4\baselineskip}\medskip\noindent\textbf{Part III --- Substitution.}\par\nopagebreak
\begin{enumerate}[resume, beginpenalty=10000]
\item State the specification of \code{fine_cents} that every policy must satisfy, as a precondition on
\code{d} and postconditions on the result, so that the fine-summing code is correct for every policy.
\item The library's reminder letter promises ``the fine never exceeds 5.00''. Which policies satisfy that promise?
Is it part of the specification of \code{FinePolicy}?
\item A client test suite checks, for every policy and every $d$ from 0 to 400, that the fine is non-negative and
non-decreasing in $d$. Do the three policies pass? Does \code{CappedFine}?
\item A clerk proposes a policy for damaged items that returns $-500$ (a refund) when $d = 0$. Which property of
question 13 does it break, and what could the summing code then compute wrongly?
\item Is \code{StudyFine} a substitute for \code{StandardFine}, in the sense of the principle, for a client that
relies on the 5.00 ceiling? For a client that relies only on question 13?
\item Should \code{ComicFine} derive from \code{StandardFine} because both cap the fine? Argue with the principle.
\end{enumerate}

\needspace{4\baselineskip}\medskip\noindent\textbf{Part IV --- The price of dispatch.}\par\nopagebreak
\begin{enumerate}[resume, beginpenalty=10000]
\item In \cref{out:b2:cpp-polymorphism:dispatch-cost}, how many more nanoseconds does a virtual call on a mix of
kinds cost than the switch?
\item The library computes fines for its 400 loans once a night. Estimate the time spent in dispatch.
\item A simulation updates $10^8$ particles of two kinds per second with one virtual call each. Estimate the share
of a second spent in calls, from the same output.
\item Name two changes to the simulation that would reduce that share, from the end of
\cref{sec:b2:cpp-polymorphism:design}.
\item Why is the measured difference between ``one kind'' and ``mixed kinds'' not explained by the two loads of
\cref{prop:b2:cpp-polymorphism:layout}?
\item State the result of the problem: what adding a kind costs in each design, in files and places edited, and what
a virtual call costs against a switch on this machine.
\end{enumerate}
\end{problem}
